% !TeX root = ../../Book3.tex
\section{次调和函数}
\label{b3:sec:1503}

全纯函数在零点处的对数模趋于负无穷，不能作为通常的调和函数处理。把平均等式放宽为不等式，并允许这个无穷值，便能保留最大值原理和零点信息。

\subsection{次调和函数}
\label{b3:subsec:150301}

\begin{definition}\label{b3:def:15-subharmonic}
设 \(D\) 为区域，\(u:D\to[-\infty,\infty)\) 上半连续、局部可积且不恒为 \(-\infty\)。若每个内部闭圆盘满足
\[
u(a)\le\frac1{2\pi}\int_0^{2\pi}u(a+r e^{\mathrm{i}t})\,\mathrm{d}t,
\]
则称 \(u\) 为\term[cidiaohe-hanshu]{次调和函数}[subharmonic function]。圆周积分取扩展意义。函数 \(-u\) 称为\term[chaodiaohe-hanshu]{超调和函数}[superharmonic function]。
\end{definition}
\(C^2\) 函数次调和当且仅当 \(\Delta u\ge0\)：必要性用二阶平均展开，充分性用圆周平均导数的面积公式。因此 \(|z|^2\) 次调和而非调和，\(-|z|^2\) 超调和。局部可积条件保证后面的分布运算有定义，也排除处处无穷的退化情形。

\subsection{平均值不等式}
\label{b3:subsec:150302}

圆盘上的比较描述同样有用。若 \(h\) 在圆盘调和、连续到边界，且 \(u\le h\) 于边界，则 \(u\le h\) 于内部。证明可先对 \(u-h+\varepsilon|z|^2\) 使用严格平均不等式：它不能有内部最大值，随后令 \(\varepsilon\downarrow0\)。一般上半连续边值用连续上逼近处理。反过来，连续边值的 Poisson 解与 \(u\) 比较，就给中心平均不等式。

\begin{proposition}\label{b3:prop:15-subharmonic-operations}
有限个次调和函数的最大值次调和；非负线性组合只要不恒为负无穷也次调和。若 \(u_j\downarrow u\) 且极限局部可积，则 \(u\) 次调和。
\end{proposition}
\begin{proof}
取最大值的函数在中心选出达到最大的一项，将其平均用所有项的最大值控制。和的情形由积分线性性。下降极限上半连续；在固定圆周上减去局部共同上界，再用单调收敛，得到平均不等式。局部可积保证其为允许的对象。
\end{proof}
任意上确界一般不保持上半连续；这正是 Perron 法必须处理而不能略去的一点。

\subsection{最大值原理}
\label{b3:subsec:150303}

\begin{proposition}[次调和函数的最大值原理]\label{b3:prop:15-subharmonic-maximum}
次调和函数若在区域内部取得有限最大值，则恒等于该值。若 \(D\) 有界且每个边界点满足 \(\limsup_{z\to\xi}u(z)\le0\)，则 \(u\le0\) 于 \(D\)。
\end{proposition}
\begin{proof}
在最大点，圆周平均与中心值相同。若圆周某处小于最大值，上半连续性给一段严格下降的弧，使平均严格下降，矛盾。故最大值集合既开又闭。边界条件加有限覆盖使某个边界邻域中 \(u\le\varepsilon\)；若内部仍有值大于 \(\varepsilon\)，最大值将在余下的紧集取得，第一结论迫使常数，和边界矛盾。令 \(\varepsilon\downarrow0\)。
\end{proof}

\subsection{\texorpdfstring{\(\log|f|\)}{log|f|}}
\label{b3:subsec:150304}

\begin{proposition}\label{b3:prop:15-log-subharmonic}
非恒零全纯函数 \(f\) 的 \(\log|f|\) 次调和，零点处取 \(-\infty\)。
\end{proposition}
\begin{proof}
局部写 \(f(z)=(z-a)^m g(z)\)，\(g\) 无零点。函数 \(\log|g|\) 调和，\(\log|z-a|\) 局部可积。后者以连续函数 \(\tfrac12\log(|z-a|^2+\varepsilon^2)\) 下降逼近；其 Laplace 算子为 \(2\varepsilon^2/(|z-a|^2+\varepsilon^2)^2\ge0\)。于是它次调和；加上调和函数保持不等式。局部性质通过圆盘比较描述给全域性质。
\end{proof}
这一证明保留零点重数，后面会把 \(m\) 解释为 Laplace 测度在 \(a\) 的质量。不能在零点任意补一个有限值；那会破坏平均不等式；仅仅保持上半连续性并不足够。


严格平均不等式的比较证明可以避免对 \(-\infty\) 点求导。设圆盘 \(B\) 上 \(u\le h\) 于边界，若内部存在 \(u-h>0\)，选充分小 \(\varepsilon>0\)，使 \(u-h+\varepsilon(|z|^2-R^2)\) 仍在某个内部点为正，而边界非正。上半连续性使它在闭圆盘取得有限正最大值。以最大点为中心取一个内部小圆，\(|z|^2\) 的圆周平均比中心多 \(r^2\)，所以中心值严格小于平均，和最大性矛盾。该证明只使用上半连续性、平均不等式及一个严格次调和修正，适用于对数奇性。

\ProseParagraphBreak

在圆盘边界上，上半连续且有上界的数据可由连续函数从上逼近。具体可先将数据截断到 \([-N,M]\)，再用小邻域上确界和平滑分片构造连续上界；截断后的单调逼近回到扩展值。Poisson 核非负且总质量一，使这种上逼近与积分相容。这个过程说明比较描述和圆周平均描述为什么一致，也解释了为什么不能用任意的点态修改代替上半连续代表。

\ProseParagraphBreak

一个常见构造是凸函数作用于调和函数。若 \(u\) 调和、\(\Phi:\mathbb R\to\mathbb R\) 凸且连续，则 Jensen 不等式与平均值性质给
\[
\Phi(u(a))\le\frac1{2\pi}\int_0^{2\pi}\Phi(u(a+re^{\mathrm{i}t}))\,\mathrm{d}t.
\]
故 \(\Phi\circ u\) 次调和。取 \(\Phi(t)=|t|\)、\(e^t\) 可得到许多不是调和的函数；在 \(C^2\) 情况，\(\Delta\Phi(u)=\Phi''(u)|\nabla u|^2\) 显示非调和性集中在凸性的二阶项。凸性不能删，例如 \(\Phi(t)=-t^2\) 与 \(u=x\) 给 Laplace 为负。

\ProseParagraphBreak

两个实部组合成复全纯函数时，\(|f|^p\) 对每个 \(p>0\) 次调和。在无零点局部，用调和函数 \(\log|f|\) 与凸函数 \(e^{pt}\) 即得；在零点处，函数连续、中心值零，平均非负，故平均不等式直接成立。这个结论不需要 \(p\ge1\)，因为这里使用的是指数作用于对数模，而不是直接声称 \(t\mapsto|t|^p\) 在所有实数上凸。
